% ============================================================
%  CENG 467 — TEXT-TO-SQL SEMANTIC PARSING
%  MID-PHASE PROGRESS REPORT
%  Izmir Institute of Technology, Spring 2026
% ============================================================

\documentclass[runningheads]{llncs}

% ---- PAKETLER ----
\usepackage[T1]{fontenc}
\usepackage[utf8]{inputenc}
\usepackage{graphicx}
\usepackage{booktabs}
\usepackage{multirow}
\usepackage{amsmath}
\usepackage{hyperref}
\usepackage{url}
\usepackage{listings}
\usepackage{xcolor}
\usepackage{enumitem}

% ---- KOD BLOĞU STİLİ ----
\lstset{
  basicstyle=\small\ttfamily,
  keywordstyle=\color{blue}\bfseries,
  commentstyle=\color{gray},
  stringstyle=\color{orange},
  frame=single,
  breaklines=true,
  language=SQL
}

% ============================================================
\begin{document}

\title{Prompt-Based Text-to-SQL Semantic Parsing\\with Schema Serialization Strategies}

\titlerunning{Prompt-Based Text-to-SQL Parsing}

\author{
  Ay\c{s}enur Sava\c{s}\inst{1} \and
  Mustafa [SOYADI]\inst{1}
}

\authorrunning{A.~Sava\c{s} and M.~[Soyad{\i}]}

\institute{
  Department of Computer Engineering,
  Izmir Institute of Technology (IYTE), Izmir, Turkey \\
  \email{\{aysenursavas, mustafa.[soyadi]\}@iyte.edu.tr}
}

\maketitle

% ============================================================
%  1. PROJECT TITLE, GROUP MEMBERS, AND PROBLEM DEFINITION
% ============================================================

\section{Project Title, Group Members, and Problem Definition}
\label{sec:problem}

\noindent\textbf{Project Title:} Prompt-Based Text-to-SQL Semantic Parsing with Schema Serialization Strategies

\noindent\textbf{Group Members:}
\begin{itemize}[noitemsep]
  \item Ay\c{s}enur Sava\c{s}
  \item Mustafa [SOYADI]
\end{itemize}

\noindent\textbf{Project Number:} 12 — Text-to-SQL Semantic Parsing

\subsection{Problem Definition}

Text-to-SQL semantic parsing is the task of automatically translating natural language questions into executable SQL queries over relational databases.
This task is central to building natural language interfaces to databases (NLIDB), enabling non-expert users to query structured data without knowledge of SQL syntax~\cite{Yu2018Spider}.

\noindent\textbf{Task formulation:}
\begin{itemize}[noitemsep]
  \item \textbf{Input:} A natural language question $q$ and a database schema $\mathcal{S}$
  \item \textbf{Output:} An executable SQL query $y$ that answers the question
  \item \textbf{Task type:} Conditional generation (sequence-to-sequence)
\end{itemize}

The task is technically challenging due to several factors:
(1)~the compositional and hierarchical nature of SQL (nested queries, joins, aggregations),
(2)~the need for accurate \textit{schema linking} between question entities and database columns/tables,
and (3)~the requirement for \textit{cross-domain generalization} where models must handle unseen database schemas at test time~\cite{Yu2018Spider}.

We use the \textbf{Spider benchmark}~\cite{Yu2018Spider} as our primary dataset and Google Gemini~\cite{GeminiTeam2024} as the language model for prompting-based baselines.
In the final submission, we plan to incorporate instruction-tuned fine-tuning and compare it against prompting strategies.

Our project title has not changed since the proposal stage.


% ============================================================
%  2. DATASET / BENCHMARK STATUS
% ============================================================

\section{Dataset / Benchmark Status}
\label{sec:dataset}

We use the \textbf{Spider} benchmark~\cite{Yu2018Spider}, a large-scale, cross-domain Text-to-SQL dataset.

\begin{table}[ht]
  \centering
  \caption{Spider dataset statistics.}
  \label{tab:dataset}
  \begin{tabular}{ll}
    \toprule
    \textbf{Property} & \textbf{Value} \\
    \midrule
    Dataset name       & Spider (xlangai/spider on HuggingFace) \\
    Data source        & HuggingFace Datasets hub \\
    Training examples  & 7{,}000 \\
    Validation examples & 1{,}034 \\
    Number of databases & 200+ across 138 domains \\
    Input format       & Natural language question + database schema \\
    Output format      & SQL query string \\
    \bottomrule
  \end{tabular}
\end{table}

\noindent\textbf{Preprocessing steps completed:}
\begin{itemize}[noitemsep]
  \item Dataset downloaded and loaded via the HuggingFace \texttt{datasets} library
  \item Database schema information extracted by parsing SQL queries across the training set (table and column names identified per database)
  \item Three schema serialization formats implemented (plain text, CREATE TABLE syntax, compact notation)
  \item Zero-shot and few-shot prompt templates constructed
\end{itemize}

\noindent\textbf{Known challenges:}
\begin{itemize}[noitemsep]
  \item The HuggingFace version of Spider does not include explicit schema metadata (\texttt{tables.json}); we reconstruct schema information by parsing SQL queries from the training set
  \item Schema linking remains imperfect for complex multi-table queries
  \item Some SQL queries involve nested subqueries and complex aggregations that are difficult to generate via prompting alone
\end{itemize}

The dataset has been \textbf{downloaded, cleaned, and prepared} for experiments.


% ============================================================
%  3. LITERATURE REVIEW PROGRESS
% ============================================================

\section{Literature Review Progress}
\label{sec:literature}

We have reviewed the following key areas relevant to our project:

\subsection{Sequence-to-Sequence Approaches for Text-to-SQL}

Early neural approaches to Text-to-SQL adopted sequence-to-sequence architectures.
Zhong et al.~\cite{Zhong2017Seq2SQL} introduced Seq2SQL, which used reinforcement learning to generate SQL from natural language over single tables.
The Spider benchmark~\cite{Yu2018Spider} later introduced cross-domain evaluation, requiring models to generalize to unseen databases.
These works established the standard task formulation and evaluation metrics (Exact Match and Execution Accuracy) that we use in our project.

\subsection{Pre-trained Language Models for Text-to-SQL}

The introduction of large pre-trained language models has substantially advanced Text-to-SQL performance.
BRIDGE~\cite{Lin2020BridgeV2} proposed a hybrid approach combining pre-trained representations with schema linking.
UnifiedSKG~\cite{Xie2022UnifiedSKG} demonstrated that multi-task learning across structured knowledge grounding tasks improves Text-to-SQL accuracy.
More recently, large language models (LLMs) such as GPT-4 and Gemini have shown strong zero-shot and few-shot performance on Text-to-SQL benchmarks~\cite{Rajkumar2022Evaluating, Gao2023DAILSQL}.
DAIL-SQL~\cite{Gao2023DAILSQL} achieved competitive results with carefully designed prompting strategies, demonstrating that prompt engineering is a viable alternative to fine-tuning for this task.

\subsection{Schema Encoding and Linking}

Effective schema encoding is critical for cross-domain generalization.
RAT-SQL~\cite{Wang2020RATSQL} introduced relation-aware attention for encoding schema structures.
PICARD~\cite{Scholak2021PICARD} augmented T5 with constrained decoding to ensure syntactic validity.
Our project focuses on a simpler approach: comparing different schema \textit{serialization} formats within prompts, which to our knowledge has not been systematically studied with instruction-tuned LLMs.

\noindent\textbf{Positioning of This Work.}
Unlike prior approaches that rely on complex graph-based schema linking, we adopt a prompting-based framework with Gemini and systematically compare schema serialization strategies. We plan to extend this with instruction tuning for the final submission.


% ============================================================
%  4. BASELINE MODELS OR PROMPTING STRATEGIES
% ============================================================

\section{Baseline Models or Prompting Strategies}
\label{sec:baselines}

We have designed and implemented two baseline prompting strategies:

\subsection{Baseline 1: Zero-Shot Prompting}

We evaluate Google Gemini (1.5 Flash)~\cite{GeminiTeam2024} in a zero-shot setting.
The prompt provides the database schema and the natural language question without any SQL examples.
This establishes the floor performance achievable by prompting alone without task-specific demonstrations.

\begin{lstlisting}[language={},caption={Zero-shot prompt template}]
You are an expert SQL developer. Your task is to
translate the natural language question into SQL.

Database: {db_id}
Schema: {serialized_schema}
Question: {question}

Return ONLY the SQL query, nothing else.
\end{lstlisting}

\subsection{Baseline 2: Few-Shot Prompting (3-shot)}

We extend the zero-shot baseline by providing three SQL examples before the target question.
Examples are selected from the same database when available, otherwise generic examples are used.
This tests whether in-context learning with examples improves SQL generation quality.

\noindent\textbf{Current status:}
\begin{itemize}[noitemsep]
  \item Both baselines are \textbf{fully implemented} in our evaluation pipeline (\texttt{evaluate.py})
  \item Schema serialization supports three formats (Format~A: plain text, Format~B: CREATE TABLE, Format~C: compact notation)
  \item API quota limitations are currently being resolved; initial test runs have been conducted
  \item Systematic evaluation on the full validation set is in progress
\end{itemize}


% ============================================================
%  5. INITIAL EXPERIMENTAL RESULTS
% ============================================================

\section{Initial Experimental Results}
\label{sec:results}

Due to API quota limitations encountered during the evaluation period, we report results from our initial test runs on a small subset.
Full evaluation on the Spider validation set is in progress and will be completed for the final submission.

\begin{table}[ht]
  \centering
  \caption{
    Preliminary results on Spider validation subset.
    EM = Exact Match (\%).
  }
  \label{tab:preliminary}
  \begin{tabular}{lccc}
    \toprule
    \textbf{Method} & \textbf{Schema} & \textbf{EM (\%)} & \textbf{Status} \\
    \midrule
    Zero-shot (Gemini 1.5 Flash)   & Format A & --    & In progress \\
    Zero-shot (Gemini 1.5 Flash)   & Format B & --    & In progress \\
    Few-shot 3-shot (Gemini 1.5 Flash) & Format A & --    & In progress \\
    Few-shot 3-shot (Gemini 1.5 Flash) & Format B & --    & In progress \\
    \midrule
    Instruction-tuned (planned)    & Best     & --    & Not started \\
    \bottomrule
  \end{tabular}
\end{table}

\noindent\textbf{Preliminary observations:}
From our initial API tests, we observed that:
\begin{itemize}[noitemsep]
  \item Zero-shot prompting produces syntactically valid SQL for simple queries but struggles with multi-table joins and nested subqueries
  \item Schema information significantly affects output quality; omitting the schema leads to frequent hallucination of non-existent table names
  \item Few-shot examples help the model follow the expected output format more consistently
\end{itemize}

We expect to obtain complete numerical results once API access is fully configured.


% ============================================================
%  6. PLANNED IMPROVEMENTS AND TECHNICAL DIRECTION
% ============================================================

\section{Planned Improvements and Technical Direction}
\label{sec:improvements}

For the next phase, we plan the following improvements:

\begin{enumerate}[noitemsep]
  \item \textbf{Complete baseline evaluation:} Run both baselines on the full Spider validation set (1{,}034 examples) with all three schema formats
  \item \textbf{Instruction-tuned fine-tuning:} Fine-tune a smaller model (Flan-T5-base~\cite{Chung2022FlanT5} or similar) with LoRA~\cite{Hu2022LoRA} on Spider training data, using the best-performing schema format
  \item \textbf{Improved schema extraction:} Download and integrate Spider's original \texttt{tables.json} for accurate schema information including primary and foreign keys
  \item \textbf{Better few-shot selection:} Implement similarity-based example selection (selecting examples whose questions are semantically similar to the target)
  \item \textbf{Post-processing:} Add SQL syntax validation and basic error correction as a post-processing step
  \item \textbf{Execution accuracy:} Implement execution-based evaluation by running generated SQL against actual databases
\end{enumerate}


% ============================================================
%  7. ABLATION OR PROMPT SENSITIVITY PLAN
% ============================================================

\section{Ablation or Prompt Sensitivity Plan}
\label{sec:ablation}

Our ablation study will systematically compare the following dimensions:

\begin{table}[ht]
  \centering
  \caption{Planned ablation study dimensions.}
  \label{tab:ablation_plan}
  \begin{tabular}{ll}
    \toprule
    \textbf{Dimension} & \textbf{Variants} \\
    \midrule
    Schema format      & Format A (plain text) vs.\ B (CREATE TABLE) vs.\ C (compact) \\
    Prompting strategy & Zero-shot vs.\ 1-shot vs.\ 3-shot vs.\ 5-shot \\
    Temperature        & 0.0 vs.\ 0.3 vs.\ 0.7 \\
    Model size         & Gemini 1.5 Flash vs.\ Gemini 1.5 Pro (if available) \\
    With/without fine-tuning & Prompting only vs.\ LoRA fine-tuned \\
    \bottomrule
  \end{tabular}
\end{table}

The primary research question is: \textit{How does schema serialization format affect Text-to-SQL accuracy across different prompting strategies?}
We hypothesize that Format~B (CREATE TABLE syntax) will perform best with code-oriented models, as it resembles the input distribution seen during pre-training.


% ============================================================
%  8. ERROR ANALYSIS AND GENERATION QUALITY ANALYSIS
% ============================================================

\section{Error Analysis and Generation Quality Analysis Plan}
\label{sec:error_analysis}

For the final report, we plan to conduct a detailed error analysis by:

\begin{enumerate}[noitemsep]
  \item \textbf{Manual inspection} of at least 50 failure cases from the validation set
  \item \textbf{Categorization} of errors into the following types:
  \begin{itemize}[noitemsep]
    \item Schema linking errors (wrong table or column selected)
    \item Incorrect JOIN conditions
    \item Aggregation errors (COUNT, SUM, AVG misuse)
    \item Nested query failures (missing or incorrect subqueries)
    \item WHERE clause mismatches (wrong conditions or values)
    \item Hallucinated table/column names
  \end{itemize}
  \item \textbf{Difficulty-level analysis:} Comparing performance across Spider's difficulty categories (easy, medium, hard, extra hard)
  \item \textbf{Generation quality:} Evaluating syntactic validity, semantic correctness, and query efficiency
\end{enumerate}

We will also examine cases where the generated SQL produces correct results but differs from the gold SQL (valid alternatives), which is captured by execution accuracy but missed by exact match.


% ============================================================
%  9. ETHICAL CONSIDERATIONS AND BIAS
% ============================================================

\section{Ethical Considerations and Bias}
\label{sec:ethics}

Deploying Text-to-SQL systems introduces several ethical considerations:

\noindent\textbf{Hallucination Risk:}
LLMs may generate syntactically valid SQL that queries non-existent columns or produces incorrect results silently.
Non-expert users cannot easily detect such errors, which could lead to incorrect data-driven decisions.

\noindent\textbf{SQL Injection Risk:}
Integrating user-provided free-text into SQL generation pipelines could introduce injection vulnerabilities if outputs are executed without sanitization.

\noindent\textbf{Domain and Language Bias:}
Spider's training data is English-only and skews toward standard business domains (universities, music, sports).
Models trained on Spider may perform poorly on multilingual queries or specialized domains such as healthcare or legal databases, raising fairness concerns.

\noindent\textbf{API Usage:}
We use the Google Gemini API for inference. All API calls use standard, publicly available endpoints, and no private or sensitive data is sent through the API.
We are mindful of responsible API usage and rate limiting.


% ============================================================
%  10. GITHUB AND REPRODUCIBILITY STATUS
% ============================================================

\section{GitHub and Reproducibility Status}
\label{sec:github}

\noindent\textbf{Repository:} \url{https://github.com/Aysenursvs/CENG467-Text-to-SQL}

\noindent\textbf{Current repository structure:}
\begin{lstlisting}[language={},caption={Project directory structure}]
CENG467-Text-to-SQL/
  src/
    data_prep.py     # Dataset download and exploration
    evaluate.py      # Baseline evaluation pipeline
    utils.py         # Schema extraction, prompts, metrics
    train.py         # (Planned) Fine-tuning script
  data/              # Dataset statistics
  results/           # Evaluation results (JSON)
  notebooks/         # Exploration notebooks
  report/            # This LaTeX report
  requirements.txt   # Python dependencies
  README.md          # Setup and usage instructions
  .env               # API configuration (not committed)
  .gitignore
\end{lstlisting}

\noindent\textbf{Reproducibility checklist:}
\begin{itemize}[noitemsep]
  \item[\checkmark] GitHub repository created and actively maintained
  \item[\checkmark] README with setup instructions and usage commands
  \item[\checkmark] \texttt{requirements.txt} with all Python dependencies
  \item[\checkmark] \texttt{.env} template for API key configuration (documented in README)
  \item[\checkmark] \texttt{.gitignore} excluding data, results, and sensitive files
  \item[\checkmark] Consistent commit history (8+ commits as of checkpoint)
  \item[$\square$] Automated evaluation script with command-line arguments
  \item[$\square$] Full results reproduction with single command
\end{itemize}


% ============================================================
%  11. CURRENT CHALLENGES AND NEXT STEPS
% ============================================================

\section{Current Challenges and Next Steps}
\label{sec:challenges}

\subsection{Current Challenges}

\begin{itemize}[noitemsep]
  \item \textbf{API quota limitations:} The free-tier quota of the Gemini API has been exhausted during development, delaying systematic evaluation. We are exploring alternative API plans and model endpoints.
  \item \textbf{Schema extraction:} The HuggingFace Spider dataset does not include explicit schema metadata; we reconstruct schema information by parsing SQL queries, which is imperfect for databases with tables not referenced in the training queries.
  \item \textbf{Computational constraints:} Fine-tuning is planned using a personal computer without a dedicated GPU, which may require using parameter-efficient methods (LoRA/QLoRA) or cloud-based resources (Google Colab).
  \item \textbf{Evaluation complexity:} Exact Match is a strict metric that may underestimate actual performance; implementing Execution Accuracy requires setting up individual SQLite databases for each Spider example.
\end{itemize}

\subsection{Next Steps}

\begin{enumerate}[noitemsep]
  \item Resolve API access and complete full baseline evaluation (Week 12)
  \item Download and integrate Spider's original schema files for improved schema encoding (Week 12)
  \item Implement and evaluate instruction-tuned fine-tuning with LoRA (Week 13)
  \item Conduct ablation study across schema formats and prompting strategies (Week 13)
  \item Perform error analysis on at least 50 failure cases (Week 14)
  \item Write final report and prepare project demonstration (Week 14--15)
\end{enumerate}


% ============================================================
%  REFERENCES
% ============================================================

\bibliographystyle{splncs04}
\bibliography{references}

\end{document}
